% Requires: \usepackage{amsmath,amssymb}

\section{Hybrid amortized configuration-conditional spectral calibration model}

\subsection{Campaign-level data, metadata, and notation}

Let the training corpus be a collection of campaigns
\[
\mathcal D=\{ \mathcal C_m \}_{m=1}^{M_{\mathrm{camp}}},
\]
where campaign $m$ is represented by
\[
\mathcal C_m = \bigl(\mathcal S_m,\;K_m,\;f_m,\;c_m,\;Y_m\bigr).
\]

Here:
\begin{itemize}
    \item $\mathcal S_m \subseteq \{1,\dots,S_{\max}\}$ is the set of sensors participating in campaign $m$, with cardinality
    \[
    M_m = |\mathcal S_m| \ge 2;
    \]
    \item $K_m \in \mathbb N$ is the number of aligned acquisitions in campaign $m$;
    \item
    \[
    f_m = (f_{m,1},\dots,f_{m,J_m}),
    \qquad
    f_{m,1} < \cdots < f_{m,J_m},
    \]
    is the common frequency grid retained after campaign preprocessing;
    \item
    \[
    c_m \in \mathcal X
    \]
    is the acquisition configuration, assumed constant within campaign $m$:
    \[
    c_m =
    \bigl(
    f_{{\mathrm c},m},\,
    B_m,\,
    \mathrm{RBW}_m,\,
    g_{{\mathrm{LNA}},m},\,
    g_{{\mathrm{VGA}},m}
    \bigr);
    \]
    \item
    \[
    Y_m=\{Y_{m,s,k,j}\}
    \]
    is the aligned PSD tensor in linear power units, indexed by
    \[
    s\in\mathcal S_m,\qquad
    k\in\{1,\dots,K_m\},\qquad
    j\in\{1,\dots,J_m\}.
    \]
\end{itemize}

Thus each campaign is internally fixed-config, but different campaigns may have different configurations $c_m$ and different frequency grids $f_m$.

Ancillary metadata such as campaign label, date range, and acquisition cadence may be stored for bookkeeping, filtering, or optional side-information. However, the physically relevant observation operator is determined by $c_m$ and the retained frequency grid $f_m$.

\subsection{Forward observation model}

For campaign $m$, sensor $s\in\mathcal S_m$, acquisition $k$, and frequency bin $j$, we posit
\begin{equation}
Y_{m,s,k,j}
=
G_{m,s,j}\,S_{m,k,j}
+
N_{m,s,j}
+
\varepsilon_{m,s,k,j},
\label{eq:hybrid-forward}
\end{equation}
where
\begin{itemize}
    \item $S_{m,k,j}\ge 0$ is the latent common spectrum for campaign $m$ at acquisition $k$ and bin $j$;
    \item $G_{m,s,j}>0$ is the multiplicative calibration gain of sensor $s$ in campaign $m$ at frequency $f_{m,j}$;
    \item $N_{m,s,j}\ge 0$ is the additive floor of sensor $s$ in campaign $m$ at frequency $f_{m,j}$;
    \item $\varepsilon_{m,s,k,j}$ is the residual term.
\end{itemize}

Conditionally on $(G,N,S)$, the working likelihood is Gaussian:
\begin{equation}
\varepsilon_{m,s,k,j}
\mid
G_{m,s,j},N_{m,s,j},S_{m,k,j}
\sim
\mathcal N\!\left(0,\sigma^2_{m,s,j}\right),
\label{eq:hybrid-gaussian}
\end{equation}
with campaign-, sensor-, and frequency-dependent residual variance $\sigma^2_{m,s,j}>0$.

The corresponding calibrated sensor-domain estimate is
\begin{equation}
X_{m,s,k,j}
=
\frac{\bigl[Y_{m,s,k,j}-N_{m,s,j}\bigr]_+}{G_{m,s,j}},
\label{eq:hybrid-corrected}
\end{equation}
where $[x]_+ = \max\{x,0\}$ or, in implementation, clipping to a small positive floor $\varepsilon_S>0$ for numerical stability.

\subsection{Hybrid amortized parameterization}

The campaign-specific calibration parameters are not fitted independently from scratch. Instead they are generated by one global trainable model with parameters $\theta$.

\subsubsection{Sensor and campaign encoders}

Each sensor identity $s\in\{1,\dots,S_{\max}\}$ is associated with a trainable embedding
\[
e_s \in \mathbb R^{d_s}.
\]

The acquisition configuration is encoded by a trainable map
\[
u_\theta : \mathcal X \to \mathbb R^{d_c},
\qquad
u_m := u_\theta(c_m).
\]

For each sensor $s\in\mathcal S_m$, define a sensor summary
\[
r_{m,s} \in \mathbb R^{d_r}.
\]
If sensor $s$ is visible to the encoder, then
\begin{equation}
r_{m,s}
=
\phi_\theta\!\left(
Y_{m,s,\cdot,\cdot},
f_m,
c_m
\right),
\label{eq:sensor-summary-visible}
\end{equation}
where $\phi_\theta$ is a sensor-wise encoder operating on the full time-frequency tensor of that sensor within campaign $m$. If sensor $s$ is masked during training, then
\begin{equation}
r_{m,s}=r_{\varnothing},
\label{eq:sensor-summary-masked}
\end{equation}
where $r_{\varnothing}\in\mathbb R^{d_r}$ is a learned mask token.

The campaign context is produced by a permutation-invariant aggregation of the visible sensor summaries:
\begin{equation}
h_m
=
\rho_\theta\!\left(
\{(e_s,r_{m,s}) : s\in\mathcal V_m\},
u_m
\right)
\in \mathbb R^{d_h},
\label{eq:campaign-context}
\end{equation}
where $\mathcal V_m\subseteq \mathcal S_m$ is the visible subset used by the encoder during training, and $\rho_\theta$ is permutation-invariant with respect to the sensor order.

\subsubsection{Frequency-basis parameterization}

To support campaigns with different frequency grids and to enforce smoothness with respect to physical frequency, the calibration curves are generated through basis expansions.

Let
\[
b^\gamma_\ell:\mathbb R\to\mathbb R,
\quad \ell=1,\dots,L_\gamma,
\]
\[
b^\eta_\ell:\mathbb R\to\mathbb R,
\quad \ell=1,\dots,L_\eta,
\]
\[
b^\nu_\ell:\mathbb R\to\mathbb R,
\quad \ell=1,\dots,L_\nu,
\]
be fixed frequency basis functions (for example, cubic B-splines defined over the relevant frequency range).

For each campaign $m$ and sensor $s\in\mathcal S_m$, the model generates coefficient vectors
\begin{align}
a^\gamma_{m,s}
&=
A^\gamma_\theta(e_s,r_{m,s},h_m,u_m)
\in\mathbb R^{L_\gamma},
\label{eq:coef-gamma}
\\
a^\eta_{m,s}
&=
A^\eta_\theta(e_s,r_{m,s},h_m,u_m)
\in\mathbb R^{L_\eta},
\label{eq:coef-eta}
\\
a^\nu_{m,s}
&=
A^\nu_\theta(e_s,r_{m,s},h_m,u_m)
\in\mathbb R^{L_\nu},
\label{eq:coef-nu}
\end{align}
where $A^\gamma_\theta$, $A^\eta_\theta$, and $A^\nu_\theta$ are trainable hypernetworks.

These coefficient vectors induce unconstrained campaign-specific curves
\begin{align}
\bar\gamma_{m,s,j}
&=
\sum_{\ell=1}^{L_\gamma}
a^\gamma_{m,s,\ell}\,
b^\gamma_\ell(f_{m,j}),
\label{eq:raw-gamma}
\\
\eta_{m,s,j}
&=
\sum_{\ell=1}^{L_\eta}
a^\eta_{m,s,\ell}\,
b^\eta_\ell(f_{m,j}),
\label{eq:raw-eta}
\\
\nu_{m,s,j}
&=
\sum_{\ell=1}^{L_\nu}
a^\nu_{m,s,\ell}\,
b^\nu_\ell(f_{m,j}).
\label{eq:raw-nu}
\end{align}

\subsubsection{Identifiability projection}

The forward model \eqref{eq:hybrid-forward} is invariant under the campaign-wise transformation
\[
G_{m,s,j}\mapsto a_{m,j}G_{m,s,j},
\qquad
S_{m,k,j}\mapsto a_{m,j}^{-1}S_{m,k,j},
\qquad
a_{m,j}>0.
\]
Therefore an identifiability constraint is required.

Let $G^{(0)}_{m,s,j}>0$ denote a nominal baseline gain. In the current project philosophy one may set
\[
G^{(0)}_{m,s,j}\equiv 1.
\]

We impose identifiability by projecting the unconstrained log-gain corrections to zero mean across the participating sensors:
\begin{equation}
\gamma_{m,s,j}
=
\bar\gamma_{m,s,j}
-
\frac{1}{M_m}
\sum_{r\in\mathcal S_m}
\bar\gamma_{m,r,j},
\qquad
s\in\mathcal S_m,\ j=1,\dots,J_m.
\label{eq:gamma-projection}
\end{equation}

Then the multiplicative gains are
\begin{equation}
G_{m,s,j}
=
G^{(0)}_{m,s,j}\exp(\gamma_{m,s,j}).
\label{eq:gain-definition}
\end{equation}

When $G^{(0)}_{m,s,j}\equiv 1$, \eqref{eq:gamma-projection} implies
\begin{equation}
\frac{1}{M_m}\sum_{s\in\mathcal S_m}\gamma_{m,s,j}=0
\quad\Longleftrightarrow\quad
\prod_{s\in\mathcal S_m} G_{m,s,j}^{1/M_m}=1,
\label{eq:geometric-mean-one}
\end{equation}
so the latent spectrum is defined on the campaign-specific network reference scale whose geometric mean gain equals one at each frequency bin.

\subsubsection{Positive floor and positive variance}

Define the softplus map
\[
\operatorname{sp}(x)=\log(1+e^x).
\]
Then
\begin{align}
N_{m,s,j}
&=
\operatorname{sp}(\eta_{m,s,j}),
\label{eq:noise-definition}
\\
\sigma^2_{m,s,j}
&=
\sigma^2_{\min}
+
\operatorname{sp}(\nu_{m,s,j}),
\label{eq:variance-definition}
\end{align}
where $\sigma^2_{\min}>0$ is a small fixed variance floor.

Thus the trainable model outputs campaign-specific calibration parameters
\[
\bigl\{G_{m,s,j},N_{m,s,j},\sigma^2_{m,s,j}\bigr\}_{s\in\mathcal S_m,\;j=1,\dots,J_m}
\]
through a single global parameter vector $\theta$ shared across all campaigns.

\subsection{Explicit latent-spectrum update}

The latent spectrum is not produced by a free black-box decoder. Instead, given the amortized campaign-specific parameters $(G,N,\sigma^2)$, it is computed analytically by weighted least squares.

Define
\begin{equation}
\omega_{m,s,j}
=
\frac{1}{\sigma^2_{m,s,j}}.
\label{eq:omega-definition}
\end{equation}

For any subset $\mathcal U_m\subseteq \mathcal S_m$, define the subset-specific latent-spectrum estimator
\begin{equation}
\widehat S_{m,k,j}^{(\mathcal U_m)}
=
\left[
\frac{
\sum_{s\in\mathcal U_m}
\omega_{m,s,j}\,G_{m,s,j}\,
\bigl(Y_{m,s,k,j}-N_{m,s,j}\bigr)
}{
\sum_{s\in\mathcal U_m}
\omega_{m,s,j}\,G_{m,s,j}^2
}
\right]_+.
\label{eq:subset-latent-estimator}
\end{equation}

At training time, to avoid trivial leakage into masked sensors, the latent spectrum is inferred from the visible subset:
\begin{equation}
\widehat S_{m,k,j}
:=
\widehat S_{m,k,j}^{(\mathcal V_m)}.
\label{eq:train-latent-estimator}
\end{equation}

At inference time on a new campaign, all available sensors are visible, so
\begin{equation}
\widehat S_{m,k,j}
:=
\widehat S_{m,k,j}^{(\mathcal S_m)}.
\label{eq:test-latent-estimator}
\end{equation}

The returned calibrated measurements are then
\begin{equation}
\widehat X_{m,s,k,j}
=
\frac{
\bigl[
Y_{m,s,k,j}-N_{m,s,j}
\bigr]_+
}{
G_{m,s,j}
}.
\label{eq:return-calibrated}
\end{equation}

Thus the complete output of the model for campaign $m$ is
\[
\Bigl(
\widehat X_m,\;
\widehat S_m,\;
G_m,\;
N_m,\;
\sigma^2_m
\Bigr).
\]

\subsection{Self-supervised masked-sensor training objective}

\subsubsection{Visible and masked subsets}

During training, each campaign sensor set is randomly partitioned into
\[
\mathcal S_m = \mathcal V_m \,\dot\cup\, \mathcal B_m,
\]
where
\begin{itemize}
    \item $\mathcal V_m$ is the visible subset used by the encoder;
    \item $\mathcal B_m$ is the masked subset used for self-supervised prediction.
\end{itemize}

The encoder receives only the visible sensors, but the model must still predict calibration parameters for all sensors.

\subsubsection{Masked negative log-likelihood}

The primary self-supervised objective is the masked negative log-likelihood
\begin{equation}
\mathcal L^{(m)}_{\mathrm{mask}}
=
\sum_{s\in\mathcal B_m}
\sum_{k=1}^{K_m}
\sum_{j=1}^{J_m}
\left[
\frac{
\bigl(
Y_{m,s,k,j}
-
G_{m,s,j}\widehat S_{m,k,j}
-
N_{m,s,j}
\bigr)^2
}{
2\sigma^2_{m,s,j}
}
+
\frac{1}{2}\log \sigma^2_{m,s,j}
\right].
\label{eq:mask-nll}
\end{equation}

Because $\widehat S_{m,k,j}$ is computed from the visible subset only, the masked sensors cannot be reconstructed by direct copying. The model must infer transferable campaign structure from the visible sensors, sensor identities, and configuration metadata.

\subsubsection{Optional visible reconstruction term}

To stabilize optimization, one may include a visible reconstruction term
\begin{equation}
\mathcal L^{(m)}_{\mathrm{vis}}
=
\sum_{s\in\mathcal V_m}
\sum_{k=1}^{K_m}
\sum_{j=1}^{J_m}
\left[
\frac{
\bigl(
Y_{m,s,k,j}
-
G_{m,s,j}\widehat S_{m,k,j}
-
N_{m,s,j}
\bigr)^2
}{
2\sigma^2_{m,s,j}
}
+
\frac{1}{2}\log \sigma^2_{m,s,j}
\right].
\label{eq:visible-nll}
\end{equation}

The visible term is optional because it is easier and may dominate training if over-weighted.

\subsection{Frequency smoothness regularization}

Since $G_{m,s,j}$, $N_{m,s,j}$, and $\sigma^2_{m,s,j}$ are functions of physical frequency, smoothness should be imposed with respect to the actual frequency grid $f_m$, which may be nonuniform.

Let
\[
\tilde h_m
=
\operatorname{median}\bigl(f_{m,j+1}-f_{m,j}\bigr)_{j=1}^{J_m-1},
\]
and define normalized coordinates
\[
x_{m,1}=0,
\qquad
x_{m,j}
=
\sum_{u=1}^{j-1}
\frac{f_{m,u+1}-f_{m,u}}{\tilde h_m},
\qquad
j=2,\dots,J_m.
\]

For $r=1,\dots,J_m-2$, set
\[
h^-_{m,r}=x_{m,r+1}-x_{m,r},
\qquad
h^+_{m,r}=x_{m,r+2}-x_{m,r+1}.
\]
Define the nonuniform three-point second-difference operator $D_{m,2}\in\mathbb R^{(J_m-2)\times J_m}$ row-wise by
\begin{equation}
(D_{m,2} z)_r
=
\frac{2}{h^-_{m,r}(h^-_{m,r}+h^+_{m,r})}z_r
-
\frac{2}{h^-_{m,r}h^+_{m,r}}z_{r+1}
+
\frac{2}{h^+_{m,r}(h^-_{m,r}+h^+_{m,r})}z_{r+2}.
\label{eq:nonuniform-second-difference}
\end{equation}

Then the smoothness penalty is
\begin{equation}
\mathcal R^{(m)}_{\mathrm{smooth}}
=
\sum_{s\in\mathcal S_m}
\left[
\lambda_\gamma
\left\|
D_{m,2}\gamma_{m,s,\cdot}
\right\|_2^2
+
\lambda_\eta
\left\|
D_{m,2}\eta_{m,s,\cdot}
\right\|_2^2
+
\lambda_\nu
\left\|
D_{m,2}\nu_{m,s,\cdot}
\right\|_2^2
\right],
\label{eq:smoothness-penalty}
\end{equation}
where
\[
\gamma_{m,s,\cdot}=(\gamma_{m,s,1},\dots,\gamma_{m,s,J_m})^\top,
\quad
\eta_{m,s,\cdot}=(\eta_{m,s,1},\dots,\eta_{m,s,J_m})^\top,
\quad
\nu_{m,s,\cdot}=(\nu_{m,s,1},\dots,\nu_{m,s,J_m})^\top.
\]

The gain regularizer acts on log-gain corrections, not on $G$ directly; likewise the additive floor regularizer acts on the softplus preimage $\eta$, not on $N$ directly.

\subsection{Optional reliable-sensor anchoring}

If a campaign-level reliable sensor $s_m^\star\in\mathcal S_m$ is available from a separate ranking diagnostic, one may encourage that sensor to remain close to the nominal gain baseline through
\begin{equation}
\mathcal R^{(m)}_{\mathrm{anchor}}
=
\sum_{j=1}^{J_m}
\gamma_{m,s_m^\star,j}^2.
\label{eq:anchor-penalty}
\end{equation}

This is a soft prior, not a hard equality constraint.

\subsection{Optional cross-configuration coefficient regularization}

The shared network weights already transfer information across configurations. If stronger configuration continuity is desired, one may add an explicit penalty on campaigns that share at least one sensor.

Let $\mathcal P$ be a set of campaign pairs $(m,m')$ such that
\[
\mathcal S_m \cap \mathcal S_{m'} \neq \varnothing.
\]
For each pair, define a configuration kernel
\begin{equation}
K(c_m,c_{m'})
=
\exp\!\left(
-\frac{
\|D_c(c_m-c_{m'})\|_2^2
}{\tau_c^2}
\right),
\label{eq:config-kernel}
\end{equation}
where $D_c$ rescales the configuration coordinates into comparable units and $\tau_c>0$ is a bandwidth parameter.

Then one may penalize coefficient discrepancies for sensors observed in both campaigns:
\begin{equation}
\mathcal R_{\mathrm{cfg}}
=
\sum_{(m,m')\in\mathcal P}
K(c_m,c_{m'})
\sum_{s\in\mathcal S_m\cap\mathcal S_{m'}}
\left[
\alpha_\gamma
\|a^\gamma_{m,s}-a^\gamma_{m',s}\|_2^2
+
\alpha_\eta
\|a^\eta_{m,s}-a^\eta_{m',s}\|_2^2
+
\alpha_\nu
\|a^\nu_{m,s}-a^\nu_{m',s}\|_2^2
\right].
\label{eq:config-regularizer}
\end{equation}

This term is optional. It encodes the prior that nearby configurations should induce nearby calibration curves for the same hardware, without forcing exact equality across configurations.

\subsection{Global training objective}

The global objective is
\begin{equation}
\mathcal L(\theta)
=
\frac{1}{M_{\mathrm{camp}}}
\sum_{m=1}^{M_{\mathrm{camp}}}
\mathbb E_{\mathcal V_m}
\Bigl[
\mathcal L^{(m)}_{\mathrm{mask}}
+
\lambda_{\mathrm{vis}}
\mathcal L^{(m)}_{\mathrm{vis}}
+
\mathcal R^{(m)}_{\mathrm{smooth}}
+
\lambda_{\mathrm{anchor}}
\mathcal R^{(m)}_{\mathrm{anchor}}
\Bigr]
+
\lambda_{\mathrm{cfg}}\mathcal R_{\mathrm{cfg}}
+
\lambda_{\mathrm{wd}}\|\theta\|_2^2.
\label{eq:global-objective}
\end{equation}
Here:
\begin{itemize}
    \item the expectation over $\mathcal V_m$ denotes random visible/masked sensor splits during stochastic training;
    \item $\lambda_{\mathrm{vis}},\lambda_{\mathrm{anchor}},\lambda_{\mathrm{cfg}},\lambda_{\mathrm{wd}}\ge 0$ are scalar hyperparameters;
    \item if a term is not used, its coefficient is set to zero.
\end{itemize}

If external reference calibrations are available for some campaigns, a supervised term may be added. For example, if reference corrected measurements $X^{\mathrm{ref}}_{m,s,k,j}$ are available, one may augment \eqref{eq:global-objective} by
\begin{equation}
\mathcal L_{\mathrm{sup}}
=
\sum_{m\in\mathcal M_{\mathrm{sup}}}
\sum_{s\in\mathcal S_m}
\sum_{k=1}^{K_m}
\sum_{j=1}^{J_m}
\left(
\widehat X_{m,s,k,j}
-
X^{\mathrm{ref}}_{m,s,k,j}
\right)^2,
\label{eq:optional-supervised-loss}
\end{equation}
with an additional multiplier $\lambda_{\mathrm{sup}}\ge 0$.

\subsection{Inference on a new campaign}

Let
\[
\mathcal C_\dagger
=
\bigl(
\mathcal S_\dagger,\,
K_\dagger,\,
f_\dagger,\,
c_\dagger,\,
Y_\dagger
\bigr)
\]
be a new campaign. At inference time all available sensors are visible:
\[
\mathcal V_\dagger = \mathcal S_\dagger.
\]
The trained model produces
\[
\{G_{\dagger,s,j},N_{\dagger,s,j},\sigma^2_{\dagger,s,j}\}_{s\in\mathcal S_\dagger,\;j=1,\dots,J_\dagger}
\]
through equations \eqref{eq:sensor-summary-visible}--\eqref{eq:variance-definition}, then computes
\[
\widehat S_{\dagger,k,j}
=
\left[
\frac{
\sum_{s\in\mathcal S_\dagger}
\omega_{\dagger,s,j}G_{\dagger,s,j}
\bigl(
Y_{\dagger,s,k,j}-N_{\dagger,s,j}
\bigr)
}{
\sum_{s\in\mathcal S_\dagger}
\omega_{\dagger,s,j}G_{\dagger,s,j}^2
}
\right]_+,
\]
and finally returns
\[
\widehat X_{\dagger,s,k,j}
=
\frac{
\bigl[
Y_{\dagger,s,k,j}-N_{\dagger,s,j}
\bigr]_+
}{
G_{\dagger,s,j}
}.
\]

Hence the trained system is a single global amortized model that maps
\[
(Y_\dagger,\mathcal S_\dagger,f_\dagger,c_\dagger)
\longmapsto
(\widehat X_\dagger,\widehat S_\dagger,G_\dagger,N_\dagger,\sigma^2_\dagger),
\]
while still preserving an explicit and interpretable calibration operator.

\subsection{Interpretation}

The model should be read as follows.

\begin{enumerate}
    \item There is one global parameter vector $\theta$ shared across all campaigns.
    \item For each new campaign, the network produces campaign-specific calibration parameters conditioned on the campaign measurements, the participating sensors, and the acquisition configuration.
    \item The latent common spectrum is then recovered by an explicit weighted least-squares consensus step rather than by an unconstrained black-box decoder.
    \item The returned calibrated PSD tensor is obtained by subtracting the inferred additive floors and dividing by the inferred multiplicative gains.
\end{enumerate}

Thus the model is simultaneously:
\begin{itemize}
    \item \emph{configuration-conditional}, because the generated calibration depends on $c_m$;
    \item \emph{campaign-specific}, because the outputs are inferred separately for each campaign;
    \item \emph{amortized}, because one shared network replaces fitting a completely separate optimization model from scratch for every campaign;
    \item \emph{hybrid and interpretable}, because the forward calibration equation remains explicit.
\end{itemize}
